
\documentclass[11pt]{article}
\usepackage[margin=1in]{geometry}
\usepackage{booktabs}
\usepackage{longtable}
\title{NPH Subgroup Scenario Simulation Study}
\author{Automated simulation pipeline}
\date{\today}
\begin{document}
\maketitle

\section*{Objective}
Following the uploaded protocol (CONFIRMS, revision 2), we implemented a simulation study for the \textbf{biomarker subgroup} non-proportional hazards scenario. The aim was to compare operating characteristics of simple analysis strategies under differential treatment effects in subgroups.

\section*{Design choices implemented}
\begin{itemize}
  \item Two-arm randomized trial, 1:1 allocation, total sample size $n=400$.
  \item Biomarker prevalence fixed at 30\%.
  \item Piecewise-constant simplification: exponential event times with arm/subgroup-specific hazards.
  \item Independent right censoring via exponential censoring model.
  \item Scenarios: null (no treatment effect), mild subgroup effect, and strong subgroup effect.
\end{itemize}

\section*{Methods evaluated}
\begin{itemize}
  \item Overall two-sample log-rank test.
  \item Subgroup-specific log-rank tests (biomarker positive / negative).
  \item Interaction-style heterogeneity test based on subgroup log rate-ratio contrast.
\end{itemize}

\section*{Parallelization and runtime control}
The simulation program uses Python multiprocessing and automatically estimates a feasible number of runs from a pilot benchmark, splitting work across all available CPU cores. The runtime budget was constrained to under 30 minutes on this machine.

\section*{Results}
\begin{center}
\begin{tabular}{lrrrrrrr}
\toprule
Scenario & $N_{sim}$ & Rej overall & Rej subgroup+ & Rej subgroup- & Rej interaction & Mean events \\
\midrule
null & 6000 & 0.057 & 0.052 & 0.051 & 0.049 & 282.4 \\
subgroup\_mild & 6000 & 0.086 & 0.172 & 0.051 & 0.142 & 279.4 \\
subgroup\_strong & 6000 & 0.152 & 0.360 & 0.053 & 0.272 & 277.7 \\
\bottomrule
\end{tabular}
\end{center}

Interpretation:
\begin{itemize}
  \item Under the null scenario, rejection rates are close to nominal 0.05 (type I error behavior).
  \item Under subgroup-effect scenarios, overall tests dilute effects, while subgroup-positive analysis gains sensitivity.
  \item Interaction rejection increases with stronger subgroup heterogeneity.
\end{itemize}

\section*{Repository structure}
\begin{verbatim}
NPH/
  configs/
  data/
  reports/
    report.tex
  results/
    summary_*.json
  scripts/
    make_report.py
  src/nph/
    simulate_subgroup.py
\end{verbatim}

\section*{Reproducibility}
Run:
\begin{verbatim}
python3 src/nph/simulate_subgroup.py --max-minutes 25 --outdir results
python3 scripts/make_report.py
\end{verbatim}

\end{document}
